\documentclass[12pt]{article}
\usepackage{amsmath, amsfonts, amssymb}
\title{Cluster Algebras}
\begin{document}
\maketitle
\section{Motivation by properties and examples}

\textbf{Example 1: } Let us take the sequence $x_0,x_1, \cdots $ defined  by the recurrence relation 
\[
x_{n+4}=\frac{x_{n+1}x_{n+3}+x_{n+2}^2}{x_n}
\]

with the initial values $x_0=x_1=x_2=x_3=1$
The first few terms of the sequence are 
\begin{align*}
    x_4=&\frac{x_1x_3+x_2^2}{x_0}=\frac{1\times1+1^2}{1}=2\\
    x_5=&\frac{x_2x_4+x_3^2}{x_1}=\frac{1\times2+1^2}{1}=3\\    
    x_6=&\frac{x_3x_5+x_4^2}{x_2}=\frac{1\times3+2^2}{1}=7\\
    x_7=&\frac{x_4x_6+x_5^2}{x_3}=\frac{2\times7 + 3^2}{1}=23\\
    x_8=&\frac{x_5x_7+x_6^2}{x_4}=\frac{3\times 23 + 7^2}{2}=59\\
    x_9=&\frac{x_6x_8+x_7^2}{x_5}=\frac{7\times59+23^2}{3}=314\\
    x_{10}=&\frac{x_7x_9+x_8^2}{x_6}=\frac{23\times314+59^2}{7}=1529
\end{align*}

Remarkably, all the terms in the sequence are integers, even though one might expect the opposite because of the complicated looking form of the recurrence relation. We will call this \textbf{Property $1'$.} That is, property $1'$ states that all the terms in the sequence, also called the Somos-4 sequence, are integers.\\

\textbf{Frieze patterns}:
The following is an example of a frieze pattern. It is an array that consists of $n-1$ rows. In particular, it is an $SL_2$-frieze pattern of order 7, that is it has 6 rows, and the top/bottom rows are all 1's. And each diamond satisfies the $SL_2$ property, that is if $\begin{pmatrix} a & b\\ c& d \end{pmatrix}$
is an $SL_2$ matrix, then $ad-bd=1$. Here each diamond would look like
$
\begin{array}{ccc}
     &a&  \\
     b&&c\\
     &d&
\end{array}
$ satisfying $ad-bc=1$
\[
\renewcommand{\arraystretch}{1.2}
\setlength{\arraycolsep}{6pt}
\begin{array}{cccccccccccccccccccc}
1^* &   & 1 &   & 1 &   & 1 &   & 1 &   & 1 &   & 1 &   & 1 &  &1&    &1\\
  & 1^*(\mathrm{ic}) &   & 2 &   & 3 &   & 1 &   & 2 &   & 3 &   & 1 &   & 1&   &1\\
 &   &  1^*(\mathrm{ic}) &   & 5 &   & 2 &   & 1 &   & 5 &   & 2 &   & 1 &     &1&  &1\\
  & 1^*(\mathrm{ic}) &   & 2 &   & 3 &   & 1 &   & 2 &   & 3 &   & 1 &   & 1&   &1\\
1^* &   & 1 &   & 1 &   & 1 &   & 1 &   & 1 &   & 1 &   & 1 &     &1&  &1\\
\end{array}
\]

\textbf{Theorem: } Any $SL_2$ frieze pattern of order $n$ is $n$-periodic.\\

\textbf{Property 2:}(Half-periodicity)\\

The initial conditions in the frieze pattern have to be chosen carefully so that the rest of the entries in the frieze pattern will be integers, For example, here the initial entries are chosen to be 1,1,1, and we get that the rest of the entries are integers.\\

\textbf{Definition:}(Cluster)  The entries of a frieze pattern form a \textit{cluster} if they lie on a lattice path, connecting the top boundary to the boundary, with steps going from top in the down-right direction and then down-left direction to the bottom. In the above example, the cluster variables are the $1's$ marked with the asterisk

\section*{Grassmanian}

Fix $0\leq k\leq n$. The \textit{Grassmanian}, denoted by $Gr(k,n)$ is defined to be 
\[
\mathrm{Gr}(k,n)=\{ V \subset \mathbf{C}^n \mid \mathrm{dim}V=k\}
\]
. In other words, the Grassmanian of $k-$planes in $\mathbb{C}^n$ is the set of all $k-$ dimensional vector subspaces in $\mathbb{C}^n$ or $\mathbb{R}^n$. For example, in $\mathbb{R}^2$, this would be the set of all lines passing through the origin, or in $\mathbb{R}^3$, this would be the set of all planes passing through the origin etc.\\

So, if $\V$ is inside $\mathrm{Gr}(k,n)$, then $V$ would have the basis $\{v_1,v_2. \cdots , v_k\}$, where $0\leq k\leq n$. If we put them in a matrix, then we would have a matrix 
\[
\begin{pmatrix}
\text{---} v_1 \text{---} \\
\text{---} v_2 \text{---} \\
\vdots \\
\text{---} v_k \text{---}
\end{pmatrix}
\]
This matrix is not unique, as choosing different bases would give us different matrices.\\

Now any matrix of this form has rank $k$, as every row is a basis vector, and spans $V$.\\
Also, this matrix would be related by left multiplication of invertible $k\times k$ matrices, that is $GL_k(\mathbb{C})$.(Since the two matrices are related simply by a change of matrices for the subspace $V)$

So, we have Gr$(k,n)=\{\text{full rank $k\times n$ matrices}\}/\{GL_k(\mathbb{C})\}$.
Therefore the dimension  dim $\mathrm{Gr}(k,n)=k(n-k)$
\subsection*{Plucker coordinates}
Let $[n]=\{1,2,\cdots,n\}$. Let $$\binom{[n]}{k}=\{J\subseteq [n]: \mid J \mid=k\}$$\\

For all $J \in \binom{[n]}{k}$, define $P_J=\text{determinant of submatrix of $M$ with column set $J$}$

\end{document}
